% =====================================================================
%  Capitolo 6 — Prospettive e struttura finale attesa
%  Nessuna conclusione: gli esperimenti non sono ancora completi.
% =====================================================================
\chapter{Prospettive e lavoro futuro}
\label{cap:prospettive}

Gli esperimenti di questa tesi non sono ancora completi, quindi il capitolo \emph{non} trae
conclusioni sulle domande di ricerca. Riassume invece lo stato del lavoro, elenca le attività
rimanenti, discute le direzioni di ricerca aperte dal quadro teorico e fissa la struttura attesa
del capitolo conclusivo.

\section{Stato del lavoro}
\label{sec:stato-lavoro}

\begin{table}[H]
  \centering
  \small
  \caption[Stato del lavoro]{Stato delle componenti del lavoro.}
  \label{tab:stato-lavoro}
  \begin{tabular}{|p{5.6cm}|l|}
    \hline
    \textbf{Componente} & \textbf{Stato} \\
    \hline
    Lettura critica della tesi del 2022 & completata (\cref{cap:premesse}) \\
    Quadro teorico & completato (\cref{cap:teoria}) \\
    Protocollo sperimentale & definito (\cref{cap:metodologia}) \\
    Implementazione e test automatici & completati, 109 test \\
    Preparazione di Potsdam, Vaihingen e pool di Toronto & completata \\
    Benchmark di memoria su RTX 4060 & completato (\cref{tab:vram}) \\
    Bootstrap per tile nel report & da implementare \\
    Campagne di esperimenti & da eseguire (\cref{tab:stato-campagne}) \\
    Analisi dei risultati e conclusioni & da scrivere \\
    \hline
  \end{tabular}
\end{table}

\section{Attività rimanenti}
\label{sec:attivita}

L'ordine di esecuzione segue la dipendenza tra le domande e il costo delle campagne:

\begin{enumerate}
  \item \codice{thesis2022} e \codice{rq1\_attention}: fissano il riferimento e quantificano
  l'effetto delle sole correzioni di protocollo (\cref{sec:exp-2022,sec:exp-rq1});
  \item \codice{thesis\_grid}: pre-addestramenti DINO R18/R34 e griglia della tesi sul protocollo
  standard, che alimenta \RQ{2} e il punto al 100\% della \RQ{3};
  \item \codice{rq6\_ndsm} e \codice{vaihingen}: relativamente brevi, verificano la
  generalizzazione su un secondo benchmark;
  \item \codice{rq3\_low\_label} e \codice{rq4\_aerial}: le campagne più lunghe, che riusano i
  pre-addestramenti già eseguiti;
  \item \codice{rq5\_efficiency} completa e \codice{transformers}, valutando se escludere
  ViT-B/16 (circa 30 minuti per epoca sulla GPU disponibile);
  \item analisi statistica (seed appaiati e bootstrap), figure qualitative, sintesi.
\end{enumerate}

\Segnaposto{Aggiornare l'elenco con le date effettive di esecuzione e con eventuali campagne
aggiunte o tolte.}

\section{Prospettive di ricerca}
\label{sec:prospettive-ricerca}

Indipendentemente dall'esito delle campagne, il quadro teorico del \cref{cap:teoria} indica
alcune direzioni che questa tesi non esplora fino in fondo.

\paragraph{Auto-supervisione densa.}
L'\cref{oss:globale-denso} osserva che la perdita di DINO agisce dopo una media globale, mentre
la segmentazione è un compito denso. Una variante con perdita locale, che confronti le
caratteristiche di regioni corrispondenti tra viste diverse, allineerebbe l'obiettivo di
pre-addestramento al compito a valle. Lo stesso vale per la ricostruzione
mascherata~\cite{he2022mae,cong2022satmae}, che opera per costruzione sulle singole porzioni
dell'immagine.

\paragraph{Diagnostica degli encoder pre-addestrati.}
Una valutazione lineare o $k$-NN sull'encoder congelato separerebbe la qualità della
rappresentazione dall'effetto del fine-tuning. L'informazione mutua~\eqref{eq:dino-mi}, stimata a
ogni epoca, darebbe un criterio di arresto più informativo dell'entropia del solo teacher.

\paragraph{Modelli di fondazione.}
Encoder pre-addestrati su collezioni molto grandi e diverse, come DINOv2~\cite{oquab2024dinov2} o
modelli specifici per il telerilevamento~\cite{wang2022ssl_rs}, sono un riferimento naturale per
la \RQ{3}: dicono se il limite è nella quantità di dati di dominio o nell'obiettivo di
pre-addestramento.

\paragraph{Quantizzazione oltre il collo di bottiglia.}
Se la \cref{prop:aggiramento} trova conferma empirica, la quantizzazione delle connessioni di
salto, o una quantizzazione gerarchica come in VQ-VAE-2~\cite{razavi2019vqvae2}, è il modo per
verificare davvero l'ipotesi di un vantaggio delle rappresentazioni discrete. I codici appresi
potrebbero anche essere interpretati come \enquote{parole visive} del dominio aereo e confrontati
con le classi.

\paragraph{Attenzione efficiente.}
L'\cref{oss:contrasto} suggerisce che il LAM aggreghi contesto globale con poca selettività.
Nuclei di somiglianza più contrastati, ma ancora fattorizzabili
come~\eqref{eq:attenzione-lineare-generale}, o attenzioni locali a finestra come in
Swin~\cite{liu2021swin}, sono alternative con costo lineare da confrontare a parità di budget.

\paragraph{Fusione multimodale e generalizzazione.}
Oltre alla fusione precoce dell'nDSM, la fusione tardiva a due rami~\cite{audebert2018beyond} e il
trasferimento tra dataset (Potsdam IR-R-G $\to$ Vaihingen, stesso falso colore) misurerebbero
quanto i modelli generalizzano a città e sensori diversi.

\paragraph{Efficienza e inferenza a bordo.}
La motivazione energetica del 2022 resta valida. Quantizzazione dei pesi, distillazione verso
encoder più piccoli~\cite{hinton2015distilling} e misure di energia per tile, oltre che di
tempo, renderebbero la \RQ{5} direttamente utile per l'inferenza su piattaforme con risorse
limitate.

\section{Struttura attesa del capitolo conclusivo}
\label{sec:struttura-conclusioni}

A esperimenti completati, questo capitolo verrà sostituito da un capitolo di conclusioni con la
struttura seguente.

\begin{enumerate}
  \item \textbf{Sintesi del lavoro}: problema, revisione del 2022, protocollo e implementazione.
  \item \textbf{Risposte alle domande di ricerca}: per ciascuna, l'esito della
  \cref{tab:exp-sintesi}, con l'entità dell'effetto, la sua variabilità e le condizioni in cui
  vale.
  \item \textbf{Confronto con la tesi del 2022}: quali conclusioni di allora reggono con il
  protocollo corretto, quali cambiano e per quale motivo.
  \item \textbf{Limiti}: una sola GPU e tre seed, due soli dataset urbani europei, pre-addestramento
  limitato a DINO, valutazione di parte dei fattori solo su Potsdam.
  \item \textbf{Sviluppi futuri}: le direzioni della \cref{sec:prospettive-ricerca} ordinate in
  base ai risultati ottenuti.
\end{enumerate}

\Segnaposto{Scrivere le conclusioni solo dopo aver completato il \cref{cap:esperimenti}.}
